% ==============================================================================
% 00_introduction.tex — UE SIS589
% ==============================================================================
\chapter*{Introduction Générale}
\addcontentsline{toc}{chapter}{Introduction Générale}
\label{chap:intro}

\epigraph{%
  \itshape « La sécurité n'est pas un produit que l'on achète~;
  c'est un processus que l'on construit. »
}{--- Bruce Schneier}

% ── Contexte ──────────────────────────────────────────────────────────────────
\section*{Contexte et enjeux}

Le paysage des menaces informatiques a radicalement évolué au cours de la dernière
décennie. Les attaques ne ciblent plus seulement les grandes entreprises
occidentales~: les PME, les administrations publiques, les opérateurs de
télécommunications et les institutions financières d'Afrique centrale sont
désormais exposés aux mêmes vecteurs — rançongiciels, compromission de
comptes privilégiés, exfiltration de données, fraude au virement.

Au Cameroun, la loi n°~2010/012 relative à la cybersécurité et à la
cybercriminalité confère à l'ANTIC des prérogatives élargies~; elle impose
aux opérateurs et aux entreprises des obligations de déclaration d'incidents
et de conservation des traces numériques. La Convention de Malabo (Union
africaine, 2014) renforce ce cadre à l'échelle continentale. Ces textes ne
sont plus seulement des contraintes formelles~: ils définissent la frontière
entre la faute professionnelle et la bonne foi prouvée par les journaux.

Dans ce contexte, la maîtrise de la supervision, de l'investigation numérique
et de la réponse aux incidents est devenue une compétence cardinale pour
tout ingénieur en sécurité des systèmes d'information.

% ── Positionnement du cours ───────────────────────────────────────────────────
\section*{Positionnement de l'UE SIS589}

\begin{boxinfo}
\textbf{Prérequis recommandés~:}
\begin{itemize}
  \item Réseaux informatiques~: TCP/IP, firewalls, VLAN, protocoles applicatifs.
  \item Administration Linux et Windows au niveau opérationnel.
  \item Cours d'Audit des SI (GES4062) et/ou de Hacking Éthique (SEC4062)
        appréciés en préalable.
  \item Notions de droit numérique~: responsabilité civile, cybercriminalité.
\end{itemize}
\end{boxinfo}

SIS589 est le maillon \textbf{opérationnel} de la chaîne de compétences en
sécurité. Là où GES4062 enseigne à identifier et qualifier les risques, et
où SEC4062 enseigne à simuler les attaques, SIS589 enseigne à
\textbf{détecter, investiguer et maîtriser} ce qui survient en temps réel.

Le cours est structuré en quatre parties~:

\begin{enumerate}
  \item \textbf{Supervision et détection~:} Architecture SOC, SIEM, journaux,
        systèmes de détection d'intrusion.
  \item \textbf{Cadre légal et normatif~:} Standards internationaux, droit
        camerounais, admissibilité des preuves numériques.
  \item \textbf{Investigation numérique (forensics)~:} Principes, artefacts
        système, mémoire vive, trafic réseau, cloud.
  \item \textbf{Réponse aux incidents et amélioration continue~:} Cycle
        ISO~27035, CSIRT, threat intelligence, RETEX.
\end{enumerate}

Trois laboratoires immersifs sur environnement simulé (Wazuh + GNS3/EVE-NG)
complètent le parcours théorique.

% ── Le cas fil rouge ──────────────────────────────────────────────────────────
\section*{Le cas fil rouge~: LiZbiZ-SIS}

\begin{boxlizbiz}
\textbf{Présentation de LiZbiZ-SIS}

LiZbiZ-SIS est l'extension sécurité de l'entreprise LiZbiZ introduite dans
les cours GES4062 et SEC4062. C'est une PME e-commerce et logistique
camerounaise de 120~employés, dont l'infrastructure comprend~:

\begin{itemize}
  \item \textbf{VLAN~10 (DMZ)~:} Serveur Web Apache, reverse proxy Nginx,
        accessible depuis Internet.
  \item \textbf{VLAN~20 (Applicatif)~:} Serveur ERP, base de données MySQL,
        serveur de messagerie.
  \item \textbf{VLAN~30 (Utilisateurs)~:} 80~postes Windows~10,
        Active Directory (Windows Server~2022).
  \item \textbf{VLAN~40 (IoT Logistique)~:} Scanners codes-barres,
        terminaux mobiles.
  \item \textbf{Réseau d'administration~:} Firewall pfSense, switch manageable,
        serveur de supervision.
\end{itemize}

LiZbiZ-SIS traite quotidiennement des données personnelles de clients
(noms, adresses, historiques de commandes, transactions mobile money).
Sa surface d'exposition est réelle~: serveur web Internet-facing,
intégration Orange Money / MTN MoMo, connexions VPN avec des partenaires
logistiques.

À l'ouverture du cours~: aucun SOC, aucun SIEM, journaux conservés
7~jours localement, équipe de 2~techniciens. Une récente alerte de la
banque partenaire sur des transactions suspectes a révélé une intrusion
remontant à plusieurs semaines~: le MTTD réel était supérieur à 21~jours.
\end{boxlizbiz}

Ce scénario initial — malheureusement banal — motive l'intégralité du
parcours pédagogique. Chaque chapitre apportera un élément de réponse
concret à la question~: \textit{comment LiZbiZ-SIS peut-elle détecter,
comprendre et maîtriser une telle compromission, aujourd'hui et à l'avenir~?}

% ── Organisation pratique ────────────────────────────────────────────────────
\section*{Organisation pratique}

\begin{table}[H]
\centering
\small
\caption{Organisation du cours SIS589}
\label{tab:orga-cours}
\rowcolors{2}{TableauPair}{TableauImpair}
\begin{tabular}{p{2.5cm}p{4cm}p{7cm}}
\toprule
\tableauHeader{Partie} & \tableauHeader{Chapitres} &
\tableauHeader{Compétences visées} \\
\midrule
Supervision & Ch.~1--3 & Architecturer un SOC, configurer un SIEM,
                          déployer des IDS/IPS \\
Cadre légal & Ch.~4 & Appliquer les normes et le droit applicable \\
Forensics & Ch.~5--7 & Acquérir, préserver et analyser des preuves
                        numériques \\
IR \& Amélioration & Ch.~8--10 & Piloter la réponse à incident, exploiter
                                   la CTI, capitaliser \\
Évaluation & Ch.~11 & Certifier les acquis \\
\midrule
Labs (3) & Lab~1--3 & SIEM, forensics disque, simulation Tabletop \\
\bottomrule
\end{tabular}
\end{table}

\begin{boxresume}
\textbf{Ce que vous saurez faire à l'issue de SIS589~:}
\begin{itemize}
  \item Déployer et opérer un SOC minimal (Wazuh + ELK) adapté à une
        organisation de taille intermédiaire.
  \item Lire, corréler et interpréter des journaux hétérogènes pour
        reconstruire la chronologie d'une attaque.
  \item Acquérir une image disque légalement admissible, en extraire les
        artefacts forensiques pertinents.
  \item Piloter la réponse à un incident selon la norme ISO~27035,
        rédiger un rapport de post-mortem.
  \item Opérer dans le cadre juridique camerounais et international,
        produire des preuves recevables en justice.
\end{itemize}
\end{boxresume}
